% ─────────────────────────────────────────────────────────────────────────────
% DETAILED METHODS SECTION
% gcr-dosimetry-pipeline — Life Sciences in Space Research submission
% ─────────────────────────────────────────────────────────────────────────────

\section{Methods}
\label{sec:methods}

All computation is implemented in Python~3.9+ as the open-source package
\code{gcr-dosimetry\discretionary{-}{}{-}pipeline}.
The pipeline is structured as a sequence of composable modules, each responsible
for one physical layer of the problem:
spectrum generation~(\S\ref{sec:spectrum}),
slab transport~(\S\ref{sec:transport}),
organ dose~(\S\ref{sec:organdose}),
cancer risk~(\S\ref{sec:reid}),
SEP acute hazard~(\S\ref{sec:sep}),
uncertainty propagation~(\S\ref{sec:uq}),
and RBE benchmarking~(\S\ref{sec:rbemethods}).
Where simplifying assumptions are made (particularly in the transport
layer), we state them explicitly and discuss their domain of validity.

% ─────────────────────────────────────────────────────────────────────────────
\subsection{Mission Trajectory Model}
\label{sec:trajectory}

A Hohmann-transfer trajectory from Earth (1.0~AU) to Mars (1.524~AU) is
generated by computing heliocentric distance as a function of mission elapsed
time, $r(t)$, using Keplerian orbital mechanics:
\begin{equation}
  r(t) = \frac{a(1-e^2)}{1 + e\cos\theta(t)},
  \label{eq:orbit}
\end{equation}
where $a = (r_\oplus + r_\mathrm{Mars})/2 = 1.262$~AU is the semi-major axis
and $e = (r_\mathrm{Mars} - r_\oplus)/(r_\mathrm{Mars} + r_\oplus) = 0.208$
is the eccentricity of the transfer ellipse.
Transit duration is fixed at 259~days (one half-period of the transfer orbit).
The heliocentric distance enters the dose rate only through the solar modulation
potential $\phi(t)$; direct geometric flux dilution as $r^{-2}$ is not applied
separately because the force-field modulation already encodes the energy
dependence of solar wind effects at the relevant energies.

For surface exposure at Mars, a constant heliocentric distance of 1.524~AU is
assumed with an Olympus Mons atmospheric overburden of approximately
$22\,\mathrm{g\,cm^{-2}}$, though the present paper focuses on the transit
phase where the SEP and shielding comparisons are most directly constrained
by MSL/RAD data.

% ─────────────────────────────────────────────────────────────────────────────
\subsection{Solar Modulation Potential}
\label{sec:phi}

The time-resolved solar modulation potential $\phi(t)$ is obtained from the
Usoskin et~al.\ (2017) database, which provides monthly $\phi$ values
reconstructed from ground-level neutron monitor count rates
back to 1936.
For the MSL transit period (2011~November -- 2012~July), values cluster around
$\phi \approx 481 \pm 50$~MV corresponding to an ascending phase toward
solar maximum.
For future mission scenarios, $\phi$ is frozen at a user-specified
value representing an assumed solar activity level.

We note a known limitation of the Usoskin database: systematic offsets between
different neutron monitor stations can introduce $\pm 50$--$100$~MV uncertainty
in the absolute $\phi$ calibration.
This uncertainty is explicitly propagated in Section~\ref{sec:uq} via a
multiplicative scale parameter on $\phi$.

% ─────────────────────────────────────────────────────────────────────────────
\subsection{GCR Spectrum Generation}
\label{sec:spectrum}

\subsubsection{Force-Field Modulation}
\label{sec:forcefield}

Solar modulation of the GCR flux is computed using the Gleeson--Axford (1968)
force-field approximation, which treats the heliosphere as a spherically
symmetric potential well of depth $\phi$ (in MV) through which particles must
propagate diffusively.
For an ion of charge $Z$ and mass number $A$, the modulated differential flux
at kinetic energy per nucleon $E$ is:
\begin{equation}
  J(E, \phi) = J_\mathrm{LIS}(E + \delta E)
  \cdot \frac{E(E + 2E_0)}{(E + \delta E)(E + \delta E + 2E_0)},
  \label{eq:ff}
\end{equation}
where $E_0 = 938.272$~MeV is the proton rest-mass energy,
$\delta E = (Z/A)\,\phi$ (MeV/nucleon) is the kinetic energy shift due to
solar deceleration, and $J_\mathrm{LIS}$ is the local interstellar spectrum
described in Section~\ref{sec:lis}.
The ratio in Eq.~\eqref{eq:ff} is the Jacobian of the energy transformation,
ensuring flux conservation.

The force-field approximation is known to be accurate for rigidities above
approximately $1$~GV ($E/n \gtrsim 200$~MeV for protons) and for solar
modulation potentials in the range $200$--$1500$~MV \citep{usoskin2017}.
At lower energies, charge-sign dependent drift effects (important for
anomalous cosmic rays and at solar minimum) are not captured.
For this application, where dose is dominated by particles above a few
hundred MeV/n, this limitation is not expected to significantly change
the results.

\subsubsection{Local Interstellar Spectra}
\label{sec:lis}

\paragraph{Protons.}
The LIS for protons follows the parameterization of Vos \& Potgieter \cite{vos2015}:
\begin{equation}
  J_\mathrm{LIS}^\mathrm{H}(E) = C_\mathrm{H}
  \cdot \frac{(E_\mathrm{GeV} + 0.67)^{-3.93}}{\beta^2},
  \label{eq:lis_proton}
\end{equation}
where $C_\mathrm{H} = 1.9 \times 10^4$~cm$^{-2}$~s$^{-1}$~MeV$^{-1}$~sr$^{-1}$,
$E_\mathrm{GeV} = E / 1000$ is the kinetic energy in GeV/nucleon, and
$\beta = v/c$ is the particle velocity normalized to the speed of light.
The $\beta^{-2}$ factor reflects the well-known ``1/v'' rise in flux at
low energies due to the slowing of particles by ionization energy loss.
This parametrization was fit to Voyager~1 and PAMELA measurements and
is considered one of the more reliable proton LIS forms currently available,
though low-energy Voyager constraints remain subject to ongoing refinement
as the spacecraft continues to recede from the heliosphere.

\paragraph{Helium.}
The helium LIS uses the same functional form with species-specific parameters
($C_\mathrm{He} = 1.82 \times 10^3$, spectral index $\alpha_\mathrm{He} = 2.78$,
offset $= 0.62$), which produces a He/H ratio of approximately $0.096$ at
$1$~GeV/n, consistent with ACE CRIS measurements \citep{stone1998}.
The harder helium spectrum ($\alpha_\mathrm{He} = 2.78$ vs.\ proton
$\alpha_\mathrm{H} = 3.93$) means helium contributes proportionally more
flux at high energies.

\paragraph{Heavy ions.}
Carbon, oxygen, silicon, and iron LIS parametrizations follow the per-species
HelMod model of Boschini et al.\ \cite{boschini2020}, which provides explicit spectral indices
and normalizations for each species fit to ACE, PAMELA, and AMS-02 data:
\begin{equation}
  J_\mathrm{LIS}^{(s)}(E_n) = C_s \cdot
  \frac{(E_\mathrm{GeV/n} + \delta_s)^{-\alpha_s}}{\beta^2}.
  \label{eq:lis_heavy}
\end{equation}
Iron is particularly important: with $\alpha_\mathrm{Fe} = 2.65$ (flatter than
protons), there is significantly more high-energy iron flux, and via
$Z^2$ scaling of stopping power, iron contributes disproportionately to
dose equivalent at all depths.
The adopted spectral parameters are summarized in Table~\ref{tab:lis_params}.

\begin{table}[H]
  \centering
  \caption{LIS spectral parameters used in \code{gcr/spectrum.py}.
           Proton and helium from Vos \& Potgieter \cite{vos2015}; heavy ions from
           Boschini et al.\ \cite{boschini2020}.}
  \label{tab:lis_params}
  \begin{tabular}{lcccl}
    \toprule
    Species & $C_s$ & $\alpha_s$ & $\delta_s$ (GeV/n) & Primary source \\
    \midrule
    H       & $1.90\times10^4$ & 3.93 & 0.67 & Vos \& Potgieter (2015) \\
    He      & $1.82\times10^3$ & 2.78 & 0.62 & Vos \& Potgieter (2015) \\
    C       & $2.50\times10^2$ & 2.72 & 0.50 & Boschini et al.\ (2020) \\
    O       & $1.80\times10^2$ & 2.69 & 0.50 & Boschini et al.\ (2020) \\
    Si      & $3.00\times10^1$ & 2.70 & 0.50 & Boschini et al.\ (2020) \\
    Fe      & $4.30\times10^0$ & 2.65 & 0.50 & Boschini et al.\ (2020) \\
    \bottomrule
  \end{tabular}
\end{table}

\paragraph{Per-species normalization.}
\label{sec:calib}
A set of empirical calibration factors $f_s$ are applied to each species
flux after modulation:
\begin{equation}
  J_s(E, \phi) \leftarrow f_s \cdot J_s(E, \phi),
  \label{eq:calib}
\end{equation}
with values $f_\mathrm{H} = 0.590$, $f_\mathrm{He} = 0.082$,
$f_\mathrm{C} = 1.980$, $f_\mathrm{O} = 0.016$,
$f_\mathrm{Si} = 0.025$, $f_\mathrm{Fe} = 0.002$.
These factors were fit by minimizing the squared difference between the
pipeline's predicted dose rate and the MSL/RAD cruise measurement
$\dot{D} = 1.84$~mGy/day at $16\,\mathrm{g\,cm^{-2}}$ aluminum,
with $\phi = 481$~MV.
We emphasize that this constitutes a \emph{calibration step}, not an
independent prediction: the pipeline is calibrated to one operational
flight measurement and then used to extrapolate to other shielding
depths, materials, launch dates, and mission types.
Accordingly, the agreement between the pipeline and MSL/RAD at the
calibration point should not be interpreted as a validation of the
underlying physical model; the independent validation evidence consists
of the shielding-depth dependence, the material ratio, and the temporal
modulation (Section~\ref{sec:validation}).

Several normalization values depart substantially from unity and merit
comment.
$f_\mathrm{He} = 0.082$ and $f_\mathrm{O} = 0.016$ suppress helium and
oxygen to below 10\% of their LIS predictions, while
$f_\mathrm{C} = 1.980$ nearly doubles the carbon flux.
These large corrections arise from a fundamental underdetermination:
six species-level free parameters are fit to a single total dose scalar.
Many combinations of $\{f_s\}$ reproduce $\dot{D} = 1.84$~mGy/day
while distributing dose very differently among species.
The physical plausibility of individual $f_s$ values can be partially
assessed against ACE/CRIS measurements of GCR composition near
1~AU \citep{stone1998}, which report He/H~$\approx 0.096$ at
$\sim\!500$~MeV/n — consistent with the unmodified LIS prediction of
$\approx 0.096$, suggesting $f_\mathrm{He}$ should be near unity rather
than $0.082$.
The factor-of-five discrepancy implies that the optimizer is
suppressing helium not because the helium LIS is wrong, but because
it is trading helium flux against carbon flux (which has a
$\sim\!36\times$ higher stopping power per nucleon and therefore
disproportionate leverage on total dose) to satisfy the single-point
constraint.
$f_\mathrm{C} > 1$ when ACE/CRIS abundance ratios suggest the
Boschini (2020) carbon LIS is already near the measured value is
most plausibly explained by the missing magnesium contribution
(Section~\ref{sec:discuss_limits}): if Mg ($Z=12$, absent from the
model) accounts for $\sim\!3\%$ of dose, the optimizer inflates
whatever high-LET species has the most flexibility — here, carbon.
The normalization factors are therefore best understood as a joint
correction for missing species, absolute LIS normalization offsets, and
CSDA dose-model systematics across the species mix — not as independent
corrections to each species' intrinsic flux.
A species-resolved calibration using the MSL/RAD per-species dose
fractions \citep{zeitlin2013} would directly constrain each $f_s$
and is a planned extension of this work.

A comparison with the NASA Space Radiation Laboratory (NSRL) GCR
simulator reference field \citep{slaba2016} provides independent
context for the per-species dose fractions the pipeline produces.
Table~\ref{tab:species_compare} shows the charge-group dose breakdown
from the pipeline at the calibration geometry against the NSRL
reference at comparable conditions.

\begin{table}[H]
  \centering
  \small
  \caption{Per-charge-group absorbed dose fractions at shielded
           geometries: pipeline output vs.\ the NSRL GCR simulator
           reference field \citep{slaba2016}.
           The two configurations differ by $4\,\mathrm{g\,cm^{-2}}$
           of aluminum and by solar conditions (ascending vs.\ minimum);
           neither difference can account for the
           $\sim\!4\times$ excess in the pipeline HZE fraction.}
  \label{tab:species_compare}
  \begin{tabular}{lccc}
    \toprule
    Particle class
      & \multicolumn{1}{p{3.2cm}}{\centering Pipeline \\ $16\gcm$ Al, $\phi=481$~MV}
      & \multicolumn{1}{p{3.2cm}}{\centering NSRL GCR reference \\ $20\gcm$ Al, solar min \citep{slaba2016}} \\
    \midrule
    $Z=1$ (protons)  & 46.7\% & 73.3\% \\
    $Z=2$ (helium)   & 16.8\% & 19.2\% \\
    $Z>2$ (HZE)      & 30.9\% &  7.6\% \\
    \bottomrule
    \multicolumn{3}{l}{\small NSRL fractions exclude $\pi$/EM cascades
      and heavy-target fragments.} \\
  \end{tabular}
\end{table}

The pipeline's proton fraction is $\sim\!27$ percentage points below
the NSRL reference, while the HZE fraction is approximately four
times higher.
Carbon ions alone account for $31.2\%$ of dose at $16\gcm$ Al in the
pipeline, yet the total HZE fraction in the NSRL reference is $7.6\%$.
This confirms, quantitatively, that the normalization factors are
producing species fractions well outside the physically expected
range — a direct consequence of the underdetermined single-point
calibration described above.
The helium fraction ($16.8\%$ vs $19.2\%$) is the one charge group
where agreement is reasonable, consistent with the unmodified
LIS He/H ratio already being close to the ACE/CRIS measurement
\citep{stone1998}.

% ─────────────────────────────────────────────────────────────────────────────
\subsection{Charged Particle Transport}
\label{sec:transport}

\subsubsection{Stopping Power and CSDA Range}
\label{sec:bethe}

Energy loss of a heavy ion traversing a material slab is computed in the
continuous-slowing-down approximation (CSDA), which treats the particle as
losing energy continuously according to its instantaneous stopping power.
The stopping power for an ion of charge $Z$ and velocity $\beta$ is modeled as:
\begin{equation}
  S(E/A,\, Z) = z_\mathrm{eff}^2(E/A,\, Z) \cdot S_p(E/A),
  \label{eq:stopping}
\end{equation}
where $S_p(E/A)$ is the proton stopping power at the same energy per nucleon,
and $z_\mathrm{eff}$ is the Barkas effective charge, which accounts for
partial charge-state stripping at energies below approximately
$1$~GeV/n \citep{barkas1956}:
\begin{equation}
  z_\mathrm{eff}(E/A,\, Z) = Z \left[1 - \exp\!\left(-125\beta Z^{-2/3}\right)\right].
  \label{eq:zeff}
\end{equation}
At energies above $\sim 1$~GeV/n, $z_\mathrm{eff} \to Z$ and the scaling
reduces to the familiar $Z^2$ form.
Proton stopping powers for aluminum, polyethylene, water, and tissue are
tabulated at 500 logarithmically spaced energy points from 1~MeV to
$10^5$~MeV using NIST PSTAR/ASTAR data, and interpolated via cubic spline.

The CSDA range is:
\begin{equation}
  R(E) = \int_0^E \frac{dE'}{S(E', Z)}\,,
  \label{eq:range}
\end{equation}
evaluated by numerical integration of the stopping power table.
The exit energy of an ion entering a slab of areal density $x\,[\mathrm{g\,cm^{-2}}]$
is obtained by inverting the range relation:
\begin{equation}
  E_\mathrm{out} = R^{-1}\!\bigl(R(E_\mathrm{in}) - x_\mathrm{eff}\bigr),
  \label{eq:eout}
\end{equation}
where $x_\mathrm{eff} = x \cdot Z^2/A$ is the effective proton-equivalent
areal density for heavy ions.
If $R(E_\mathrm{in}) < x_\mathrm{eff}$, the particle is stopped ($E_\mathrm{out} = 0$).

\subsubsection{Nuclear Interaction Attenuation}
\label{sec:nuclear}

In addition to ionization energy loss, heavy ions undergo inelastic nuclear
interactions with the shielding material.
These are modeled via an exponential fluence attenuation:
\begin{equation}
  \Phi_s(E, x) = \Phi_s^{(0)}(E) \cdot \exp\!\left(-\frac{x}{\lambda_s}\right),
  \label{eq:attenuation}
\end{equation}
where $\lambda_s$ is a species-dependent nuclear interaction mean free path.
Values are taken from Wilson et~al.\ (1991) for aluminum; for polyethylene,
$\lambda_s$ is scaled by the material density ratio and effective nuclear charge.
This exponential approximation does not account for the secondary fragment
spectrum produced in nuclear fragmentation reactions.
For iron and silicon, this approximation becomes increasingly crude above
$\sim 30\,\mathrm{g\,cm^{-2}}$, where fragmentation produces significant
lighter-ion secondaries (primarily C and O) that are themselves biologically
relevant.
A more accurate treatment would require a full nuclear reaction network
(as implemented in HZETRN or Geant4/TOPAS), and this represents the
primary physical limitation of the current transport model.

The full transport step for species $s$ through shielding $x$ is:
\begin{align}
  \Phi_s^\mathrm{out}(E) &=
    \mathcal{T}_s\bigl[\Phi_s^\mathrm{in}(E), x\bigr] \notag \\
  &= \Phi_s^\mathrm{in}(E_\mathrm{in}(E, x)) \cdot
     \exp\!\left(-\frac{x}{\lambda_s}\right) \cdot
     \frac{E_\mathrm{in}(E_\mathrm{in} + 2E_0)}{E(E + 2E_0)},
  \label{eq:transport}
\end{align}
where the final ratio is a flux density Jacobian correcting for the change
in energy bin width after deceleration.

\subsubsection{Secondary Neutron Production}
\label{sec:neutrons}

Spallation reactions between GCR ions and shielding nuclei produce secondary
neutrons that contribute meaningfully to dose equivalent, particularly in
aluminum where the cross-sections are relatively high.
Rather than tracking individual neutron histories (which would require a
full Monte Carlo or deterministic neutron transport code), neutron
dose equivalent is estimated from a pre-computed look-up table of the form:
\begin{equation}
  \dot{H}_n(x, \phi, \mathrm{mat}) = \dot{H}_n^\mathrm{Al}(\phi) \cdot
  f_\mathrm{mat} \cdot g(x),
  \label{eq:neutron}
\end{equation}
where $\dot{H}_n^\mathrm{Al}(\phi)$ is the neutron dose equivalent rate for
aluminum at a reference thickness, $f_\mathrm{mat}$ is a material-dependent
yield factor (polyethylene $\approx 0.72\times$ aluminum due to lower spallation
cross-sections), and $g(x)$ is an interpolated depth-dependence from the
HZETRN benchmark dataset of Slaba et al.\ \cite{slaba2014}.
The neutron quality factor is taken as $Q_n = 10$ (ICRP-60), which is a
weighted average over the neutron energy spectrum in shielding geometries.

This tabular approach introduces approximately
$\pm 25$--$30\%$ uncertainty in the neutron contribution to dose equivalent,
arising from uncertainties in the underlying HZETRN neutron spectra and
from simplifying the material and geometry dependence.
This uncertainty is explicitly included as a free parameter in the
LHS ensemble (Section~\ref{sec:uq}).

% ─────────────────────────────────────────────────────────────────────────────
\subsection{Organ Dose Model}
\label{sec:organdose}

\subsubsection{Organ Tissue Depths}
\label{sec:depths}

Organ-specific absorbed dose and dose equivalent are computed by transporting
the modulated, shielded GCR flux through an additional layer of tissue
corresponding to the depth of each organ in the human body.
Tissue depths are taken from ICRP Publication~89 \citep{icrp89} and are
listed in Table~\ref{tab:organ_depths}.
The tissue layer is treated as water-equivalent for the purpose of stopping
power calculations.

The total areal density seen by a particle reaching organ $k$ is:
\begin{equation}
  x_\mathrm{total}^{(k)} = x_\mathrm{wall} + d_k,
  \label{eq:total_depth}
\end{equation}
where $x_\mathrm{wall}$ is the spacecraft wall areal density and $d_k$ is
the organ tissue depth in $\mathrm{g\,cm^{-2}}$.

\begin{table}[H]
  \centering
  \caption{Organ tissue depths and ICRP-60 tissue weighting factors used in
           the pipeline. Depths from ICRP Publication~89 \citep{icrp89}.}
  \label{tab:organ_depths}
  \begin{tabular}{lccc}
    \toprule
    Organ & Depth (g/cm$^2$) & $w_T$ (ICRP-60) & Cancer type modeled \\
    \midrule
    Stomach    & 8.0  & 0.12 & Gastric carcinoma \\
    Colon      & 9.5  & 0.12 & Colorectal carcinoma \\
    Lung       & 5.0  & 0.12 & Lung carcinoma \\
    Leukemia   & 4.0  & 0.12 & Myeloid leukemia \\
    Esophagus  & 4.0  & 0.05 & Esophageal carcinoma \\
    Liver      & 6.0  & 0.04 & Hepatocellular carcinoma \\
    Bladder    & 10.0 & 0.04 & Transitional cell carcinoma \\
    Thyroid    & 2.0  & 0.04 & Thyroid carcinoma \\
    Breast     & 1.5  & 0.05 & Breast adenocarcinoma \\
    Ovary      & 12.0 & 0.20 & Ovarian carcinoma \\
    Other      & 7.0  & 0.05 & Composite remainder \\
    \bottomrule
  \end{tabular}
\end{table}

\subsubsection{Absorbed Dose Rate}
\label{sec:doserate}

The absorbed dose rate in tissue at organ depth $d_k$ is:
\begin{equation}
  \dot{D}_k = \sum_s \int_{E_\mathrm{min}}^{E_\mathrm{max}}
    \Phi_s(E; x_\mathrm{total}^{(k)},\, \phi)
    \cdot S_s(E) \cdot \rho^{-1} \cdot C
    \,dE,
  \label{eq:doserate}
\end{equation}
where $\Phi_s$ is the transported differential flux of species $s$
(cm$^{-2}$~s$^{-1}$~MeV$^{-1}$~sr$^{-1}$),
$S_s(E)$ is the stopping power (MeV~cm$^2$~g$^{-1}$),
$\rho^{-1}$ converts from material to tissue units,
and $C = 1.602 \times 10^{-10}$ Gy per (MeV g$^{-1}$ cm$^{-2}$) is a unit
conversion factor.
The energy grid spans $10^1$--$10^5$~MeV/n at 200 logarithmically spaced points,
which was found to be sufficient for $<0.5\%$ convergence of the integrated dose rate.

The factor of $4\pi$ from the isotropic flux assumption
($J \, [\mathrm{sr}^{-1}] \to \Phi = 4\pi J$) is included in the
proportionality constant.
This is a standard approximation for GCR fields, which are highly isotropic
at interplanetary distances \citep{simpson1983}.

\subsubsection{Dose Equivalent Rate}
\label{sec:doseequiv}

Dose equivalent at organ $k$ is:
\begin{equation}
  \dot{H}_k = \int Q(L)\, \dot{D}_k(L) \,dL
            = \sum_s Q_s \cdot \dot{D}_{k,s},
  \label{eq:doseequiv}
\end{equation}
where $Q(L)$ is the ICRP-60 quality factor as a function of unrestricted
linear energy transfer $L$ in water:
\begin{equation}
  Q(L) = \begin{cases}
    1                & L < 10~\mathrm{keV\,\mu m^{-1}} \\
    0.32\,L - 2.2   & 10 \leq L \leq 100~\mathrm{keV\,\mu m^{-1}} \\
    300 / \sqrt{L}  & L > 100~\mathrm{keV\,\mu m^{-1}}
  \end{cases}
  \label{eq:quality}
\end{equation}
\citep{icrp60}.
For each species $s$, the dose-averaged LET $\overline{L}_D$ is computed at
the transported mean energy, and $Q_s = Q(\overline{L}_{D,s})$ is applied.
This per-species approximation (rather than integrating $Q(L)$ over the
full LET spectrum for each species) is a simplification that introduces
an estimated error of $<5\%$ for protons and helium (where the LET
distribution is relatively narrow) but could be larger for iron at
low energies where the LET-range product is highly variable.
The effective quality factor of the full GCR field is:
\begin{equation}
  Q_\mathrm{eff} = \frac{\sum_{k,s} H_{k,s}}{\sum_{k,s} D_{k,s}},
  \label{eq:qeff}
\end{equation}
which for the MSL transit geometry at $\phi = 481$~MV and $x = 16\,\mathrm{g\,cm^{-2}}$
yields $Q_\mathrm{eff} = 2.60$, within $0.8\%$ of the MSL/RAD measurement
of $Q_\mathrm{eff} = 2.62 \pm 0.14$ \citep{zeitlin2013}.
We note that this agreement is partly a consequence of the calibration
procedure and should not be taken as an independent confirmation.

% ─────────────────────────────────────────────────────────────────────────────
\subsection{Cancer Risk Model (REID)}
\label{sec:reid}

\subsubsection{Overview of the NASA Risk Framework}
\label{sec:reidoverview}

Cancer mortality risk is quantified as the Radiation Exposure in Incidence of
Death (REID), defined as the probability that a crew member will die of
radiation-induced cancer as a consequence of their mission exposure.
The pipeline implements the Cucinotta et al.\ \cite{cucinotta2013} organ-specific risk model
(NASA/TP-2013-217375), which was NASA's primary operational risk framework
through the early 2020s and remains the most widely cited basis for
published mission REID estimates.
We note that NASA has subsequently released updated risk frameworks
(e.g., the 2023 age-dependent career limit), and the model assumptions
regarding high-LET radiobiology (particularly the DDREF and the
extrapolation from mouse and atomic-bomb survivor data to astronaut
demographics) carry substantial
uncertainty that cannot be resolved with the present approach.
These limitations are discussed in Section~\ref{sec:discuss_biology}.

\subsubsection{Excess Risk Model}
\label{sec:ear}

The REID integral is:
\begin{equation}
  \mathrm{REID}(a_e, s) = \int_{a_e + L}^{a_\mathrm{LT}}
    \frac{\mu_c(a, s) \cdot S(a \mid a_e, s)}{S(a_e \mid s)} \, da,
  \label{eq:reid}
\end{equation}
where $a_e$ is the age at exposure, $L = 10$~years is the minimum latency
before radiation-induced cancer can manifest, $a_\mathrm{LT}$ is the
maximum lifespan in the life table, $\mu_c(a, s)$ is the cause-specific
cancer excess mortality rate, and $S(a \mid a_e, s)$ is the conditional
survival probability from age $a_e$ to age $a$ drawn from 2020 CDC US
life tables for sex $s$.

The excess cancer mortality rate combines excess absolute risk (EAR) and
excess relative risk (ERR) terms:
\begin{equation}
  \mu_c(a) = \lambda_c^\mathrm{background}(a) \cdot \mathrm{ERR}(D, a_e)
           + \mathrm{EAR}(D, a_e, a),
  \label{eq:excess}
\end{equation}
where $\lambda_c^\mathrm{background}$ is the age-dependent baseline cancer
mortality rate (from SEER registry data for US males/females),
and EAR and ERR are organ-specific functions of absorbed dose $D$,
age at exposure $a_e$, and attained age $a$, taken from
Tables~2--3 of Cucinotta et al.\ \cite{cucinotta2013}.

The DDREF is applied to reduce the risk at low dose rates:
\begin{equation}
  D_\mathrm{eff} = D / \mathrm{DDREF},
  \label{eq:ddref}
\end{equation}
with DDREF nominally $1.5$ for low-LET and $1.0$ for high-LET radiation.
The handling of the DDREF for mixed GCR fields (which span a wide LET range)
remains actively debated \citep{durant2013} and represents one of the
dominant sources of uncertainty in the REID calculation.

\subsubsection{Uncertainty in REID}
\label{sec:reidunc}

The REID model is applied in both point-estimate and uncertainty-weighted forms.
For the point estimate, nominal parameter values are used.
For the uncertainty ensemble (Section~\ref{sec:uq}), multiplicative scale
factors are applied to the ERR, EAR, and DDREF within the
LHS framework.
The resulting REID distribution is skewed: because REID appears in
an exponential ($1 - e^{-\mu}$), large values of the radiobiological
parameters produce disproportionately large REID values, leading to
a heavy right tail in the p95 band.

% ─────────────────────────────────────────────────────────────────────────────
\subsection{Solar Energetic Particle Acute Risk}
\label{sec:sep}

\subsubsection{Band-Function Spectrum}
\label{sec:band}

Solar energetic particle event fluence spectra are parameterized using the
Band function \citep{tylka2006}, a piecewise power law with a smooth
transition at a break energy $E_b$:
\begin{equation}
  J_\mathrm{SEP}(E) = C \times \begin{cases}
    E^{\gamma_\ell} \exp(-E/E_b) & E \leq E_\mathrm{trans}, \\
    \left[(\gamma_\ell - \gamma_h) E_b\right]^{\gamma_\ell - \gamma_h}
    \exp(\gamma_h - \gamma_\ell) \cdot E^{\gamma_h} & E > E_\mathrm{trans},
  \end{cases}
  \label{eq:band}
\end{equation}
where $E_\mathrm{trans} = (\gamma_\ell - \gamma_h) E_b$ is the transition energy,
$\gamma_\ell$ and $\gamma_h$ are the low- and high-energy spectral indices,
and $C$ is a normalization constant.

Three canonical events are modeled with parameters from Tylka \& Dietrich \cite{tylka2006}
and references therein (Table~\ref{tab:sep_events}).
The August~1972 event is included as a historical worst-case;
it predates modern space dosimetry and its spectrum is reconstructed
from ground-level neutron monitor observations and is therefore subject
to larger uncertainty than the 2003 and 2005 events.

\begin{table}[H]
  \centering
  \caption{Band-function spectral parameters for three canonical SEP events.}
  \label{tab:sep_events}
  \begin{tabular}{lcccccc}
    \toprule
    Event & $C$ & $\gamma_\ell$ & $\gamma_h$ & $E_b$ (MeV) & Duration (days) \\
    \midrule
    Aug 1972 & $4.0\times10^4$ & $-1.78$ & $-3.98$ & 32  & 2 \\
    Oct 2003 & $1.5\times10^4$ & $-1.40$ & $-3.20$ & 50  & 3 \\
    Jan 2005 & $3.0\times10^2$ & $-0.90$ & $-2.80$ & 110 & 1 \\
    \bottomrule
  \end{tabular}
\end{table}

\subsubsection{BFO Absorbed Dose}
\label{sec:bfodose}

The blood-forming organ (BFO) absorbed dose from an SEP event is computed
by transporting the event-integrated proton fluence $\Phi(E)$ (cm$^{-2}$)
through the spacecraft wall and a nominal BFO tissue depth
of $10\,\mathrm{g\,cm^{-2}}$ using CSDA transport:
\begin{equation}
  D_\mathrm{BFO} = \int_{E_\mathrm{min}}^{E_\mathrm{max}}
    \Phi(E) \cdot S_p(E_\mathrm{out}(E)) \cdot C \, dE,
  \label{eq:sep_dose}
\end{equation}
where $E_\mathrm{out}(E)$ is the proton exit energy after traversing
the combined wall~+~tissue slab, and particles that stop within this
slab deposit their residual range.
Results are compared against the NASA 30-day BFO limit of $250$~mGy
\citep{cucinotta2010}.

We treat only protons in the SEP module; heavy-ion SEP fluence at
these energies is generally negligible for dose calculations,
though it may be non-negligible for quality factor estimates at
very thin shielding.

\subsubsection{Mission Encounter Probability}
\label{sec:seprisk}

The probability of encountering at least one large SEP event during a mission
of duration $T$ days is modeled as a Poisson process:
\begin{equation}
  P_\mathrm{SEP}(T) = 1 - \exp(-\lambda T / 365.25),
  \label{eq:seprisk}
\end{equation}
where $\lambda = 0.8$~events/year is the adopted rate of large
($>10^9$~proton~cm$^{-2}$ at $>10$~MeV) events during solar maximum.
During solar minimum, $\lambda$ drops by approximately a factor of 3--5
\citep{feynman1993}, making launch timing a relevant design consideration.

% ─────────────────────────────────────────────────────────────────────────────
\subsection{Uncertainty Quantification via Latin Hypercube Sampling}
\label{sec:uq}

\subsubsection{Parameter Space}
\label{sec:params}

A global physics and radiobiological uncertainty analysis is performed using
Latin Hypercube Sampling (LHS), a stratified random sampling method that
provides better coverage of the joint parameter space than simple Monte Carlo
at the same sample count.
LHS divides each parameter's distribution into $N$ equally probable strata
and draws exactly one sample from each stratum, then randomly permutes
the samples across parameters to avoid correlation.

Eight parameters are varied simultaneously (Table~\ref{tab:uq_params}).
Four govern physics inputs (solar modulation, LIS normalization, nuclear
cross-sections, and neutron secondary yield) and four govern radiobiology
(quality factor scaling, DDREF, ERR scale, EAR scale).

\begin{table}[H]
  \centering
  \caption{Latin Hypercube Sampling parameters, distributions, and physical
           justifications. Uniform[$a$,$b$] = uniform distribution on $[a,b]$;
           LogNormal[$\mu$,$\sigma$] = lognormal with geometric mean $\mu$
           and multiplicative standard deviation $\sigma$.}
  \label{tab:uq_params}
  \begin{tabular}{llll}
    \toprule
    Parameter & Distribution & Justification \\
    \midrule
    \code{phi\_scale}    & Uniform[0.80, 1.20] & Usoskin NM inter-station offset ($\pm20\%$) \\
    \code{LIS\_norm}     & Uniform[0.85, 1.15] & Voyager/AMS-02 calibration spread \\
    \code{cross\_sec}    & Uniform[0.90, 1.10] & Nuclear $\sigma$ tables: $\pm10\%$ \\
    \code{neutron\_H}    & Uniform[0.70, 1.30] & HZETRN neutron table: $\pm30\%$ \citep{slaba2014} \\
    \code{Q\_factor}     & Uniform[0.80, 1.20] & ICRP-60 $Q(L)$ form uncertainty \\
    \code{DDREF}         & LogNormal[1.0, 0.35] & BEIR VII geometric mean + spread \\
    \code{ERR\_scale}    & LogNormal[1.0, 0.30] & Cucinotta (2013) ERR fit residuals \\
    \code{EAR\_scale}    & LogNormal[1.0, 0.30] & Cucinotta (2013) EAR fit residuals \\
    \bottomrule
  \end{tabular}
\end{table}

The lognormal distributions for the three radiobiological parameters
(DDREF, ERR, EAR) reflect the multiplicative nature of these uncertainties
and the fact that they cannot be negative.
They are parameterized such that the 95th percentile is approximately
$1.8\times$ the median, consistent with the stated uncertainty bounds
in Cucinotta et al.\ \cite{cucinotta2013}.

\subsubsection{Ensemble Execution}
\label{sec:ensemble}

For each LHS sample, a full mission dose integration is performed:
the GCR spectrum is generated at each trajectory day, transported through
the specified shielding, organ doses and quality factors are computed,
and REID is estimated.
For publication-quality results, $N = 500$ samples were used (runtime
approximately 90 minutes on a 2020 MacBook Pro with no parallelization).
For automated testing, $N = 20$ samples are used (controlled via the
\code{UNCERTAINTY\_SAMPLES} environment variable), sufficient to verify
code correctness without burdening the CI pipeline.

\subsubsection{Sensitivity Decomposition}
\label{sec:sobol}

Parameter sensitivity is estimated using rank-based Spearman correlation
$\rho_i$ between each input parameter $x_i$ and the output quantity of interest
(REID):
\begin{equation}
  S_i = \frac{\rho_i^2}{\sum_j \rho_j^2}\,.
  \label{eq:sensitivity}
\end{equation}
This normalized $\rho^2$ decomposition is an approximation to first-order
Sobol indices; it underestimates interaction effects but is computationally
straightforward and provides useful rank-ordering of parameter importance
at modest sample counts.
A proper variance-based Sobol analysis would require $\sim 10^4$ samples
\citep{saltelli2010} and is beyond the scope of this paper.

% ─────────────────────────────────────────────────────────────────────────────
\subsection{Dose-Averaged LET and Per-Species Decomposition}
\label{sec:letmethod}

The dose-averaged LET for a single ion species $s$ is:
\begin{equation}
  \overline{L}_{D,s} = \frac{\int \Phi_s(E) \cdot S_s(E)^2 \, dE}
                            {\int \Phi_s(E) \cdot S_s(E) \, dE},
  \label{eq:letd_species}
\end{equation}
where the numerator weights LET by dose (LET $\times$ flux), and the
denominator is the dose rate.
The field dose-averaged LET is then:
\begin{equation}
  \overline{L}_D = \sum_s f_s \cdot \overline{L}_{D,s}\,,
  \label{eq:letd_field}
\end{equation}
where $f_s = D_s / \sum_{s'} D_{s'}$ is the fractional dose contribution
of species $s$.
This species decomposition allows us to trace the origin of $Q_\mathrm{eff}$
to specific ion contributions, which is the key link between GCR dosimetry
and proton therapy RBE (discussed in Section~\ref{sec:let}).

% ─────────────────────────────────────────────────────────────────────────────
\subsection{Proton Therapy RBE Benchmark}
\label{sec:rbemethods}

\subsubsection{Wedenberg Model}
\label{sec:wedenberg}

The Wedenberg et~al.\ (2013) LET-dependent RBE model is:
\begin{equation}
  \mathrm{RBE}_\mathrm{Wed}(D, \overline{L}_D, \alpha_x/\beta_x) =
    \frac{\sqrt{\tfrac{1}{4}(\alpha_x/\beta_x)^2 +
          \alpha_r/\beta_x \cdot D + D^2} - \tfrac{1}{2}(\alpha_x/\beta_x)}{D},
  \label{eq:wedenberg}
\end{equation}
where $\alpha_r = \alpha_x \left(1 + q \overline{L}_D / (\alpha_x/\beta_x)\right)$,
$q = 0.434$~$\mu$m/keV is an empirical constant fitted to in vitro V79 cell data,
$\alpha_x$ and $\beta_x$ are the photon radiosensitivity parameters of the
reference cell line, and $D$ is the absorbed dose.
The model reduces to RBE~$= 1$ when $\overline{L}_D \to 0$ and
increases monotonically with LET.

\subsubsection{McNamara Model}
\label{sec:mcnamara}

The McNamara et~al.\ (2015) model uses a min-max RBE parametrization:
\begin{equation}
  \mathrm{RBE}_\mathrm{Mc}(D, \overline{L}_D, \alpha_x/\beta_x) =
    \frac{\sqrt{(\alpha_x/\beta_x)^2 + 4\,(\alpha_x/\beta_x)\,
          \mathrm{RBE}_\mathrm{max}\,D + 4\,\mathrm{RBE}_\mathrm{min}^2\,D^2}
          - \alpha_x/\beta_x}{2D},
  \label{eq:mcnamara}
\end{equation}
with:
\begin{align}
  \mathrm{RBE}_\mathrm{max} &= p_1 + p_2\,(\alpha_x/\beta_x)^{-1}\,
                                \overline{L}_D, \label{eq:rbemax}\\
  \mathrm{RBE}_\mathrm{min} &= p_3 + p_4\,(\alpha_x/\beta_x)^{1/2}\,
                                \overline{L}_D, \label{eq:rbemin}
\end{align}
where $(p_1, p_2, p_3, p_4) = (0.999064, 0.35605, 1.1012, -0.0038703)$
are from McNamara et al.\ \cite{mcnamara2015}.
Both models are implemented to match the OpenTOPAS-RBE scorer
conventions exactly, enabling direct numerical comparison.

\subsubsection{TOPAS Reference Data}
\label{sec:topasref}

TOPAS-nBio reference data for a $100$~MeV proton beam in a water phantom
(200 depth bins, $0.1$~cm each, spanning $0$--$20$~cm) were generated
using the OpenTOPAS-RBE extension for cell lines V79
($\alpha/\beta = 1.41$~Gy), prostate ($\alpha/\beta = 3.0$~Gy),
and head \& neck ($\alpha/\beta = 10.0$~Gy).
For V79, the reference dataset was taken directly from the
OpenTOPAS-RBE validation suite; for the clinical cell lines,
reference values were generated by propagating the V79 LETd spectrum
through the Wedenberg and McNamara models with the appropriate
$\alpha/\beta$ values.
Accordingly, the agreement for clinical cell lines reported in
Table~\ref{tab:rbe} reflects numerical self-consistency of the pipeline's
implementation, not an independent experimental validation.
Physical LETd values from the TOPAS-nBio scorer are used as input
to both models, so the benchmark primarily tests model implementation
fidelity rather than transport accuracy.
Readers should note this distinction when interpreting the reported
mean absolute errors.
